\id{МРНТИ 28.01.75}{}

\begin{header}
\swa{Р.К.~Ускенбаева, Абдул~Разак, А.Б.~Касымова, Ж.Б.~Кальпеева, А.М.~Анартаева}{ИНТЕЛЛЕКТУАЛЬНАЯ СИСТЕМА ГРУППОВЫХ ПОКУПОК В КОНТЕКСТЕ ЭКОНОМИКИ СОВМЕСТНОГО ПОТРЕБЛЕНИЯ: АРХИТЕКТУРА, БИЗНЕС-МОДЕЛЬ И МЕЖДУНАРОДНЫЙ ОПЫТ}

\tsp{1}Р.К.~Ускенбаева\alink{https://orcid.org/0000-0002-8499-2101},
\tsp{1}Абдул~Разак\alink{https://orcid.org/0000-0003-0409-3526},
\tsp{1}А.Б.~Касымова\alink{https://orcid.org/0000-0003-2999-5745}\envelope,
\tsp{1}Ж.Б.~Кальпеева\alink{https://orcid.org/0000-0002-4970-3095},
\tsp{2}А.М.~Анартаева\alink{https://orcid.org/0009-0001-1281-0284}
\end{header}

\begin{affil}
\tsp{1}НАО Казахский национальный исследовательский технический университет им. К.И. Сатпаева, Алматы, Казахстан,

\tsp{2}Astana IT University, Астана, Казахстан

\corrauthor{Корреспондент-автор: a.kassymova@satbayev.university}
\end{affil}

В условиях цифровизации экономики и расширения моделей экономики
совместного потребления платформы групповых покупок формируются как
самостоятельный сегмент электронной коммерции, в котором экономический
результат сделки определяется коллективным поведением пользователей.
Несмотря на рост числа исследований, посвящённых архитектуре и
бизнес-моделям таких платформ, во многих работах параметры групповой
сделки рассматриваются описательно, без формализованного количественного
анализа условий её успешного завершения.

В данной статье предлагается формализованный подход к анализу механизма
групповой покупки на основе стохастического моделирования динамики
присоединения пользователей. Целью исследования является разработка
модели групповой сделки и проведение вычислительного эксперимента,
позволяющего количественно оценить влияние ключевых параметров - размера
скидки, минимального числа участников и временного окна набора - на
вероятность успешного завершения сделки.

Методология исследования основана на моделировании процесса
присоединения пользователей в виде пуассоновского потока с
интенсивностью, зависящей от ценового стимула. На основе полученного
аналитического выражения вычисляется вероятность достижения порогового
числа участников за заданное время. Проведён вычислительный эксперимент
с варьированием параметров модели, в результате которого получены
таблицы и графики, отражающие нелинейный характер зависимости
вероятности успеха от параметров сделки и наличие эффекта критической
массы.

Результаты показывают, что параметры групповой сделки не являются
взаимозаменяемыми: размер скидки выступает инструментом краткосрочной
настройки, тогда как порог участников и временное окно формируют
структурные ограничения жизнеспособности сделки. Полученные
количественные зависимости позволяют обосновывать параметры групповых
покупок на этапе проектирования цифровых платформ без привлечения
реальных транзакционных данных.

На основе результатов моделирования предложена архитектурная
интерпретация цифровой платформы групповых покупок с аналитическим
модулем поддержки принятия решений, обеспечивающим оценку вероятности
успешного завершения сделки и формирование рекомендаций по параметрам
механизма. Представленный подход может быть использован при разработке
интеллектуальных платформ коллективного потребления, в том числе в
условиях развивающихся цифровых рынков.

{\bfseries Ключевые слова:} групповые покупки, экономика совместного
потребления, цифровая платформа, социальная коммерция, цифровая
платформа, вероятность успешности сделки, поддержка принятия решений.

\begin{header}
ОРТА ТҰТЫНУ ЭКОНОМИКАСЫ ЖАҒДАЙЫНДАҒЫ ТОПТЫҚ САУДАЛАРДЫҢ ЗИЯЛДЫ ЖҮЙЕСІ: АРХИТЕКТУРА, БИЗНЕС МОДЕЛЬ ЖӘНЕ ХАЛЫҚАРАЛЫҚ ТӘЖІРИБЕ

\tsp{1}Р.К.~Ускенбаева,
\tsp{1}Абдул~Разак,
\tsp{1}А.Б.~Касымова\envelope,
\tsp{1}Ж.Б.~Кальпеева,
\tsp{2}А.М.~Анартаева
\end{header}

\begin{affil}
\tsp{1}Қ.И. Сәтбаев атындағы Қазақ ұлттық зерттеу техникалық университеті, Алматы, Қазақстан,

\tsp{2}Astana IT University, Астана, Қазақстан,

e-mail: a.kassymova@satbayev.university
\end{affil}

Экономиканың цифрландырылуы және ортақ пайдалану экономикасы
модельдерінің кеңеюімен топтық сатып алу платформалары электрондық
коммерцияның тәуелсіз сегменті ретінде пайда болуда, мұнда транзакцияның
экономикалық нәтижесі пайдаланушылардың ұжымдық мінез-құлқымен
анықталады. Мұндай платформалардың архитектурасы мен бизнес модельдеріне
арналған зерттеулер санының артуына қарамастан, көптеген зерттеулер
топтық транзакция параметрлерін сипаттамалық түрде, оны сәтті аяқтау
шарттарының ресми сандық талдауынсыз зерттейді.

Бұл мақалада пайдаланушыларды тарту динамикасын стохастикалық
модельдеуге негізделген топтық сатып алу механизмін талдаудың формальды
тәсілі ұсынылады. Зерттеудің мақсаты - топтық транзакция моделін әзірлеу
және негізгі параметрлердің - жеңілдік мөлшерінің, қатысушылардың ең аз
санының және жалдау уақытының терезесінің - транзакцияны сәтті аяқтау
ықтималдығына әсерін сандық бағалау үшін есептеу экспериментін жүргізу.

Зерттеу әдістемесі пайдаланушыларды тарту процесін баға ынталандыруына
тәуелді қарқындылығы бар Пуассон ағыны ретінде модельдеуге негізделген.
Алынған аналитикалық өрнек белгілі бір уақыт ішінде қатысушылардың шекті
санына жету ықтималдығын есептеу үшін қолданылады. Әртүрлі модель
параметрлерімен есептеу эксперименті жүргізілді, нәтижесінде табыс
ықтималдығы мен транзакция параметрлері арасындағы байланыстың сызықтық
емес сипатын, сондай-ақ маңызды массалық әсердің болуын көрсететін
кестелер мен графиктер пайда болды.

Нәтижелер топтық транзакция параметрлерінің бір-бірімен
алмастырылмайтынын көрсетеді: же\-ңілдік мөлшері қысқа мерзімді түзету
құралы ретінде қызмет етеді, ал қатысушы шегі мен уақыт аралығы
транзакцияның өміршеңдігіне құрылымдық шектеулер жасайды. Алынған сандық
қатынастар нақты транзакция деректеріне сүйенбей, сандық платформаны
жобалау кезеңінде топтық сатып алу параметрлерін негіздеуге мүмкіндік
береді.

Модельдеу нәтижелеріне сүйене отырып, транзакцияның сәтті аяқталу
ықтималдығын бағалауды және механизм параметрлеріне негізделген
ұсыныстар жасауды қамтамасыз ететін аналитикалық шешім қабылдауды қолдау
модулі бар сандық топтық сатып алу платформасының архитектуралық
түсіндірмесі ұсынылады. Ұсынылған тәсілді ұжымдық тұтынуға арналған
интеллектуалды платформаларды әзірлеуде, соның ішінде дамып келе жатқан
сандық нарықтарда қолдануға болады.

{\bfseries Түйін сөздер:} топтық сатып алу, ортақ пайдалану экономикасы,
цифрлық платформа, әлеуметтік коммерция, цифрлық платформа,
транзакциялардың сәттілік деңгейі, шешім қабылдауды қолдау.

\begin{header}
INTELLIGENT SYSTEM OF GROUP SHOPPING IN THE CONTEXT OF THE SHARED CONSUMPTION ECONOMY: ARCHITECTURE, BUSINESS MODEL AND INTERNATIONAL EXPERIENCE

\tsp{1}R.K.~Uskenbayeva,
\tsp{1}Abdul
Razaque,
\tsp{1}A.B.~Kassymova\envelope,
\tsp{1}Zh.B.~Kalpeyeva,
\tsp{2}A.M.~Anartayeva
\end{header}

\begin{affil}
Kazakh National Research Technical University named after K.I. Satbayev, Almaty, Kazakhstan,

\tsp{2}Astana IT University, Almaty, Kazakhstan,

e-mail: a.kassymova@satbayev.university
\end{affil}

With the digitalization of the economy and the expansion of sharing
economy models, group buying platforms are emerging as an independent
segment of e-commerce, where the economic outcome of a transaction is
determined by the collective behavior of users. Despite the growing
number of studies devoted to the architecture and business models of
such platforms, many studies examine group transaction parameters
descriptively, without a formal quantitative analysis of the conditions
for its successful comple\-tion.

This article proposes a formalized approach to analyzing the group
buying mechanism based on stochas\-tic modeling of user acquisition
dynamics. The aim of the study is to develop a group transaction model
and conduct a computational experiment to quantitatively evaluate the
impact of key parameters-discount size, minimum number of participants,
and recruitment time window-on the probability of successful transaction
completion.

The research methodology is based on modeling the user acquisition
process as a Poisson flow with intensity dependent on the price
incentive. The resulting analytical expression is used to calculate the
probability of reaching the threshold number of participants within a
given time. A computational experi\-ment was conducted with varying model
parameters, resulting in tables and graphs reflecting the nonlinear
nature of the relationship between the probability of success and
transaction parameters, as well as the presence of a critical mass
effect.

The results demonstrate that group transaction parameters are not
interchangeable: the discount size serves as a short-term adjustment
tool, while the participant threshold and time window form structural
constraints on the viability of the transaction. The resulting
quantitative relationships allow for the substan\-tiation of group
purchasing parameters at the digital platform design stage without
relying on actual transaction data.

Based on the modeling results, an architectural interpretation of a
digital group purchasing platform with an analytical decision support
module is proposed, providing an assessment of the probability of
successful transaction completion and generating recommendations based
on the mechanism parameters. The presented approach can be used in the
development of intelligent platforms for collective consumption,
including in emerging digital markets.

{\bfseries Keywords:} group buying, sharing economy, digital platform,
social commerce, digital platform, trans\-action success rate, decision
support.

\begin{multicols}{2}
{\bfseries Введение.} Цифровые платформы коллективного потребления и
социальной коммерции в последние годы трансформируются в устойчивый
сегмент электронной коммерции, в котором поведение потребителей
определяется не только ценой, но и социальными взаимодействиями,
алгоритмическими рекомендациями и механизмами координации спроса.
Современные исследования подчёркивают, что ключевыми факторами
устойчивости таких платформ становятся доверие, сетевые эффекты,
архитектура взаимодействия участников и способность платформы управлять
динамикой пользовательской активности на основе данных {[}1-3{]}.

Групповые покупки (group buying) представляют собой особый тип
платформенного механизма, при котором экономический результат сделки
зависит от коллективного достижения заранее заданного условия -
минимального числа участников - в пределах ограниченного временного
окна. В отличие от традиционных моделей электронной торговли, здесь
возникает эффект «критической массы»: вероятность завершения сделки
становится нелинейной функцией параметров механизма, таких как размер
скидки, порог участников и длительность набора {[}4,5{]}. Это делает
задачи проектирования и оптимизации платформ групповых покупок
принципиально количественными и требует формализованных моделей,
способных учитывать стохастический характер пользовательских потоков.

Современная литература по социальной коммерции и платформенной экономике
демонстрирует растущий интерес к алгоритмическому управлению спросом,
включая динамическое ценообразование, персонализацию условий участия и
прогнозирование успешности сделок {[}6-8{]}. Методы аналитики больших
данных и машинного обучения позволяют выявлять закономерности поведения
пользователей, однако в ряде работ отмечается, что без формализованной
постановки условий групповой сделки такие подходы остаются слабо
воспроизводимыми и трудно переносимыми между различными рынками и
платформами {[}9{]}.

Особое значение приобретают поведенческие эффекты, связанные с риском
несостоявшихся групповых сделок. Недавние исследования показывают, что
неудачное завершение акции (недостижение порога) может снижать
готовность пользователей к повторному участию, усиливать ценовую
чувствительность и изменять оптимальные параметры последующих
предложений {[}10{]}. Это обстоятельство усиливает необходимость
количественного анализа параметров групповой сделки ещё на этапе
проектирования механизма.

Актуальность формализованных моделей возрастает для развивающихся и
переходных экономик, где цифровые рынки характеризуются высокой
неоднородностью спроса и быстрыми структурными изменениями. Республика
Казахстан относится к числу таких рынков: по данным Бюро национальной
статистики, электронная коммерция демонстрирует устойчивый рост, при
этом значительная доля транзакций осуществляется через платформенные
формы (маркетплейсы и агрегаторы) {[}11{]}. Академические исследования,
посвящённые развитию e-commerce в Казахстане, фиксируют усиление роли
онлайн-каналов в розничной торговле и высокую чувствительность
потребителей к ценовым стимулам, что делает механизмы пороговых и
скидочных сделок особенно релевантными {[}12,13{]}.

Дополнительным фактором, влияющим на архитектуру и параметры цифровых
платформ в Казахстане, является развитие трансграничной электронной
торговли в рамках Евразийского экономического союза. Исследования
показывают, что особенности таможенного регулирования, порогов
беспошлинного ввоза и логистической инфраструктуры формируют
специфические условия для платформ, ориентированных на коллективные и
ценочувствительные модели потребления {[}14{]}. В этом контексте
групповые покупки могут рассматриваться как инструмент снижения
транзакционных издержек и адаптации к регуляторным ограничениям, что
усиливает практическую значимость их формализованного анализа.

Несмотря на рост числа публикаций, посвящённых архитектуре и
бизнес-моделям цифровых платформ, сохраняется научный разрыв между
обзорно-описательными исследованиями и работами, в которых параметры
групповой сделки анализируются как формализуемая система со
стохастической динамикой присоединения пользователей и измеряемыми
выходными характеристиками - такими как вероятность успешного завершения
сделки и чувствительность результата к управляемым параметрам {[}15{]}.
В этой связи целью настоящего исследования является разработка
формализованной модели групповой сделки для цифровых платформ
коллективных покупок и проведение вычислительного эксперимента,
направленного на анализ влияния размера скидки, минимального числа
участников и временного окна набора на вероятность успешного завершения
сделки. Полученные результаты ориентированы на использование в задачах
предварительного проектирования интеллектуальных платформ групповых
покупок, включая условия развивающихся цифровых рынков.

{\bfseries Материалы и методы.} Исследование выполнено в рамках
вычислительного подхода, ориентированного на формализацию механизма
групповой сделки и количественный анализ влияния её ключевых параметров
на вероятность успешного завершения. В качестве основного инструмента
используется стохастическое моделирование процесса присоединения
пользователей к групповой покупке с последующим численным анализом
чувствительности модели.

Методология исследования включает три взаимосвязанных этапа:

1. формализацию условий успешности групповой сделки;

2. построение стохастической модели динамики присоединения пользователей;

3. проведение вычислительного эксперимента с варьированием ключевых
параметров и анализ полученных зависимостей.

\emph{Формализованная постановка задачи групповой сделки.}

Рассматривается групповая покупка, характеризуемая набором параметров:

\begin{equation}
D = \left\{ P,d,N,T \right\}
\end{equation}

где \(P\) - базовая цена товара (или услуги), \(d \in (0,1)\) -
относительный размер скидки, \(N\) - минимальное число участников,
необходимое для активации сделки, \(T\) - максимальное допустимое время
набора группы.

Сделка считается успешной, если число присоединившихся пользователей
\(U(T)\) достигает или превышает порог \(N\) за время не более \(T\):

\begin{equation}
U(T) \geq N
\end{equation}

\emph{Стохастическая модель динамики присоединения пользователей.}

Процесс присоединения пользователей к групповой покупке моделируется как
стохастический поток событий. В базовой постановке предполагается, что
количество новых участников за малый интервал времени подчиняется
пуассоновскому процессу с интенсивностью \(\lambda(d)\), зависящей от
размера скидки:

\begin{equation}
U(T) Poisson(\lambda(d)T)
\end{equation}

Зависимость интенсивности потока от скидки аппроксимируется линейной
функцией:

\begin{equation}
\lambda(d) = \lambda_{0}(1 + \alpha d)
\end{equation}

где, \(\lambda_{0}\) - базовая интенсивность присоединения при
отсутствии скидки, \(\alpha > 0\) - коэффициент чувствительности спроса
к скидке.

Такая форма зависимости используется для обеспечения интерпретируемости
модели и позволяет исследовать влияние ценового стимула на скорость
формирования группы без избыточного усложнения вычислительного
эксперимента.

\emph{Вероятность успешного завершения сделки.}

Вероятность успешного завершения групповой сделки определяется как:
\end{multicols}
\vspace{-1.5em}
\begin{equation}
P_{success}(N,d,T) = P\left( U(T) \geq N \right) = 1 - \sum_{k = 0}^{N - 1}{\frac{{(\lambda(d)T)}^{k}}{k!}e^{- \lambda(d)T}}
\end{equation}
\vspace{-1.5em}
\begin{multicols}{2}
Данное выражение позволяет напрямую оценивать влияние параметров
\(N,\ d\ и\ T\) на вероятность достижения порогового условия и
используется в качестве основной выходной метрики вычислительного
эксперимента.

\emph{Постановка вычислительного эксперимента.}

Вычислительный эксперимент проводится путём варьирования параметров
модели в заданных диапазонах:

- размер скидки \(d \in \lbrack 0.05,\ 0.10,\ 0.15,\ 0.20\rbrack\);

- минимальный порог участников \(N \in \left\{ 20,30,40 \right\}\);

- временное окно набора \(T \in \lbrack 24\rbrack\) часов.

Базовая интенсивность \(\lambda_{0}\) и коэффициент \(\alpha\)
выбираются так, чтобы отражать реалистичные сценарии пользовательской
активности на цифровых платформах групповых покупок. Для каждого набора
параметров вычисляется вероятность успешного завершения сделки согласно
(5), после чего анализируется чувствительность результата к изменению
отдельных параметров при прочих равных условиях.

\emph{Анализ чувствительности.}

Анализ чувствительности выполняется по каждому параметру отдельно:

- анализ влияния скидки \(d\) при фиксированных \(N\) и \(T\);

- анализ влияния порога \(N\) при фиксированных \(d\) и \(T\);

- анализ влияния временного окна \(T\) при фиксированных \(d\) и \(N\).

Полученные зависимости интерпретируются с точки зрения экономической
логики платформенных механизмов и используются для выявления областей
параметров, в которых вероятность успешного завершения сделки возрастает
нелинейно (эффект «критической массы»).

\emph{Ограничения методологии.}Отметим, что в рамках данной работы
используется упрощённая стохастическая модель, не учитывающая
гетерогенность пользователей, сетевые эффекты социальных связей и
динамическое изменение параметров сделки во времени. Тем не менее,
предложенная постановка позволяет получить воспроизводимые
количественные оценки и служит основой для последующего расширения
модели и её эмпирической валидации на реальных транзакционных данных при
разработке программного прототипа платформы.

{\bfseries Результаты и обсуждение.} Параметры вычислительного
эксперимента. Численные результаты получены на основе аналитического
выражения для вероятности успешного завершения групповой сделки (формула
(5)) при фиксированных параметрах базовой интенсивности
пользовательского потока и чувствительности к скидке:

\begin{equation*}
\lambda_{0} = 2,0\frac{пользователя}{час}, \alpha = 5
\end{equation*}

Данные значения соответствуют умеренному уровню пользовательской
активности, характерному для региональных цифровых платформ групповых
покупок, и используются как базовый сценарий для анализа
чувствительности.

Влияние размера скидки и порогового числа участников. В таблице 1
представлены значения вероятности успешного завершения сделки при
различных уровнях скидки и пороговых значениях числа участников при
фиксированном временном окне \(T = 24\ часа\).
\end{multicols}

\tcap{Таблица 1 - Вероятность успешного завершения групповой сделки
при различных параметрах (\(\mathbf{T = 24}\) часа)}
\begin{longtblr}[
  label = none,
  entry = none,
]{
  width = \linewidth,
  colspec = {Q[100]Q[200]Q[122]Q[122]Q[122]},
  cells = {c},
  cells = {font = \small},
  row{1} = {font=\bfseries\small},
  hlines,
  vlines,
}
Размер скидки & Ожидаемое число участников & \(\mathbf{P}_{\mathbf{success}}\mathbf{(N = 20)}\) & \(\mathbf{P}_{\mathbf{success}}\mathbf{(N = 30)}\) & \(\mathbf{P}_{\mathbf{success}}\mathbf{(N = 40)}\) \\
5\% & 36 & 0,97 & 0,79 & 0,33 \\
10\% & 48 & 0,99 & 0,96 & 0,72 \\
15\% & 60 & 0,999 & 0,995 & 0,93 \\
20\% & 72 & 0,999 & 0,999 & 0,99 
\end{longtblr}

\begin{multicols}{2}
Полученные данные показывают, что влияние скидки на вероятность успеха
носит выраженный нелинейный характер. При умеренных значениях порога
(\(N = 30\)) увеличение скидки с 5\% до 10\% приводит к росту
вероятности успешного завершения сделки на 17 процентных пунктов, тогда
как дальнейшее увеличение скидки даёт существенно меньший предельный
эффект. Это свидетельствует о наличии зоны убывающей отдачи от ценового
стимулирования.

Данная закономерность наглядно представлена на рисунке 1, где показана
зависимость вероятности успешного завершения сделки от размера скидки
при различных значениях порога участников.

\fig[\columnwidth]{i/image70}[Рис.1 - Зависимость вероятности успешного завершения сделки от
размера скидки при различных порогах $\mathbf{N(T = 24 \text{ часа})}$]

Следует отметить, что для \(N = 20\) кривая практически совпадает с
верхней границей графика, поскольку вероятность успешного завершения
сделки уже при минимальной скидке превышает 0,95, что отражает высокую
устойчивость сделки при низком пороге участников.

Роль временного окна набора группы. Временной фактор оказывает не менее
значимое влияние на успешность групповой сделки. В таблице 2 приведены
результаты для фиксированной скидки \(d = 10\%\) при различных значениях
временного окна.
\end{multicols}

\tcap{Таблица 2 - Влияние длительности временного окна на вероятность
успешного завершения сделки $\mathbf{(d = 10\%)}$}
\begin{longtblr}[
  label = none,
  entry = none,
]{
  cells = {c},
  cells = {font = \small},
  cell{1}{1} = {font=\bfseries\small},
  hlines,
  vlines,
}
Время набора, ч & $\mathbf{P_{success}(N = 20)}$      & $\mathbf{P_{success}(N = 30)}$     \\
6               & 0,35  & 0,01 \\
12              & 1,00  & 0,86 \\
24              & 0,99  & 0,96 \\
36              & 0,999 & 0,99 
\end{longtblr}

\begin{multicols}{2}
Результаты показывают, что сокращение временного окна до 6 -12 часов
резко снижает вероятность достижения порогового числа участников,
особенно при более высоких значениях \(N\). Даже при достаточно
привлекательной скидке такие акции характеризуются высокой вероятностью
неуспеха.

Графическая интерпретация данной зависимости представлена на рисунке 2,
где отчётливо виден пороговый характер влияния времени набора.

\fig[\columnwidth]{i/image71}[Рис.2 - Зависимость вероятности успешного завершения сделки от
времени набора группы при скидке 10\%]
\end{multicols}

\tcap{Таблица 3 - Минимальная скидка, обеспечивающая $\mathbf{P_{success} \ge 0,9 \ (T = 24\text{\bfseries ч})}$}
\begin{longtblr}[
  label = none,
  entry = none,
]{
  cells = {c},
  cells = {font = \small},
  row{1} = {font=\bfseries\small},
  hlines,
  vlines,
}
Порог $\mathbf{N}$ & Минимальная скидка \\
20    & 5\%                \\
30    & 10\%               \\
40    & 15\%               
\end{longtblr}

\begin{multicols}{2}
Результаты таблицы 3 могут быть напрямую использованы как правило
первичной настройки параметров сделки при проектировании платформы: рост
порога на 10 участников требует увеличения скидки примерно на 5\% для
сохранения высокой вероятности успеха.

Совместный анализ таблиц 1-3 и рисунков 1-2 показывает, что параметры
групповой сделки формируют ограниченную область допустимых значений,
внутри которой достигается высокая вероятность успешного завершения. При
выходе за границы этой области система переходит в режим,
характеризующийся резким ростом доли несостоявшихся сделок.

Особенно важно, что параметры скидки, порога и времени не являются
взаимозаменяемыми. Скидка выступает управляемым параметром краткосрочной
настройки, тогда как пороговое число участников и временное окно задают
структурные ограничения механизма. Этот эффект критической массы
подчёркивает необходимость количественного обоснования параметров сделок
при проектировании цифровых платформ.

Полученные численные зависимости имеют прямую архитектурную
интерпретацию. Они указывают на необходимость выделения в архитектуре
платформы отдельного аналитического компонента, отвечающего за оценку
вероятности успешного завершения сделки и динамическую корректировку её
параметров.

В рамках предлагаемой архитектуры такой компонент целесообразно
реализовывать в виде сервиса принятия решений, интегрированного с
модулем управления сделками. Этот сервис использует входные параметры
(текущая скорость присоединения пользователей, уровень скидки,
оставшееся время) для пересчёта вероятности успеха и инициирует события
корректировки условий акции.

На рисунке 3 представлена архитектурная схема цифровой платформы
групповых покупок. Центральным элементом системы является модуль
управления групповыми сделками, обеспечивающий координацию
взаимодействия пользователей, аналитического модуля и платёжного шлюза.
Аналитический модуль используется как компонент поддержки принятия
решений и выполняет оценку вероятности успешного завершения сделки на
основе текущей динамики набора участников, формируя рекомендации по
корректировке параметров сделки. На основании принятых решений модуль
управления инициирует финансовые операции и формирование уведомлений для
пользователей, что обеспечивает разделение вычислительной логики,
бизнес-процессов и пользовательских коммуникаций.
\end{multicols}

\fig[0.6\textwidth]{i/image72}[Рис.3 - Архитектура цифровой платформы групповых покупок с
аналитическим модулем поддержки принятия решений]

\begin{multicols}{2}
Результаты вычислительного эксперимента подтверждают, что успешность
групповой сделки определяется нелинейным взаимодействием скидки,
порогового числа участников и временного окна. Представленные таблицы и
графики формируют воспроизводимую количественную основу для обоснования
параметров платформ групповых покупок и устраняют недостатки чисто
описательных архитектурных подходов.

{\bfseries Выводы.} В работе предложена формализованная модель групповой
сделки для цифровых платформ коллективных покупок, позволяющая
количественно оценивать вероятность её успешного завершения в
зависимости от ключевых параметров --- размера скидки, минимального
числа участников и временного окна набора. Модель основана на
стохастическом описании динамики присоединения пользователей и
обеспечивает аналитическое выражение для оценки успешности сделки.

Результаты вычислительного эксперимента показали нелинейный характер
влияния параметров групповой сделки и наличие эффекта критической массы,
при котором незначительное изменение параметров приводит к резкому росту
или снижению вероятности успеха. Установлено, что размер скидки является
управляемым параметром краткосрочной настройки, тогда как пороговое
число участников и временное окно задают структурные ограничения
жизнеспособности сделки.

Полученные численные зависимости подтверждают необходимость
предварительного количественного анализа параметров групповых покупок на
этапе проектирования цифровых платформ. Представленные таблицы и графики
формируют воспроизводимую основу для обоснования управленческих решений
без привлечения реальных транзакционных данных и разработки
полнофункционального программного прототипа.

Архитектурная интерпретация результатов демонстрирует возможность
интеграции предложенной модели в цифровую платформу в виде
аналитического модуля поддержки принятия решений. Такой подход
обеспечивает разделение вычислительной логики, бизнес-процессов и
пользовательских коммуникаций и повышает обоснованность настройки
параметров групповых сделок.

В целом результаты исследования подтверждают применимость
вычислительного моделирования как инструмента анализа и проектирования
механизмов групповых покупок в условиях развивающихся цифровых рынков и
могут быть использованы при создании интеллектуальных платформ
коллективного потребления.

Следующие исследования целесообразно направить на учёт гетерогенности
пользователей, сетевых эффектов социальных связей и динамического
изменения параметров сделки во времени, а также на эмпирическую
валидацию модели на реальных данных цифровых платформ.

\emph{{\bfseries Финансирование.}} \emph{Данное исследование финансируется
Комитетом науки Министерства науки и высшего образования Республики
Казахстан (грант AP23489233 - «SmartBuy Connect: интеллектуальная
система групповых покупок на основе AI»).}
\end{multicols}

\begin{center}
{\bfseries Литература}
\end{center}

\begin{refs}
1. Duan C. A State-of-the-Art Review of Sharing Economy Business Models
and a Forecast of Future Research Directions for Sustainable
Development: A Bibliometric Analysis Approach. // Sustainability/- 2023.
- Vol.15(5). - P.4568. DOI 10.3390/su15054568.

2. Faraji M., Seifdar M.H., Amiri B. Sharing economy for sustainability:
A review// Journal of Cleaner Production. - 2024. - Vol.434: 140065.
DOI 10.1016/j.jclepro.2023.140065.

3. Tan T.M., et al. Trust in blockchain-enabled exchanges: Future
directions in platform-based markets// Journal of the Academy of
Marketing Science. - 2023. - Vol.51(4). - P.914-939. DOI
10.1007/s11747-022-00889-0.

4. Xu S., Chen T. Research on Optimal Group-Purchase Threshold and
Pricing Strategy of Community Group Purchase//Mathematics. - 2023. -
Vol.11(24):4951. DOI 10.3390/math11244951.

5. Jiang Y., Wu M., Li X. Cooperation mechanism, product quality, and
consumer segment in group buying// Electronic Commerce Research and
Applications. - 2024. - Vol.64: 101357. DOI
10.1016/j.elerap.2024.101357.

6. Kopalle P.K., et al. Dynamic pricing: Definition, implications for
managers, and future research directions// Journal of Retailing. - 2023.
- Vol.99(4). - P.580--593. DOI 10.1016/j.jretai.2023.11.003.

7. Rohn D., Bican P., Brem A., Kraus S., Clauß T. Digital platform-based
business models. An exploration of critical success factors// Journal of
Engineering and Technology Management. - 2021. - Vol.60(3). DOI
10.1016/j.jengtecman.2021.101625.

8. Herzallah D., Liébana-Cabanillas F., Muñoz-Leiva F. Et al. Past,
present, and future of social commerce: a bibliometric analysis// Qual
Quant. - 2025. - Vol.59. - P.3079--3111. DOI
10.1007/s11135-025-02103-z.

9. Xu S, Chen T. Research on Optimal Group-Purchase Threshold and
Pricing Strategy of Community Group Purchase//Mathematics.
-2023.-Vol.11(24):4951. DOI 10.3390/math11244951.

10. Wang J., et al. Group buying with consumer disappointment at failed
deals// Transportation Research Part E. - 2025. - Vol.194: 103903. DOI
10.1016/j.tre.2024.103903.

11. Bureau of National Statistics of the Republic of Kazakhstan.
E-commerce in the Republic of Kazakhstan (2024). {[}Электронный
ресурс{]}.2025. - Режим доступа:
\url{https://stat.gov.kz/en/industries/economy/local-market/publications/370122/}
- Дата обращения: 10.02.2026.

12. Жанбозова А.Б., Азатбек Т.А., Дадабаева Д.М., Кашкинбаева К. Бөлшек
сауда сегментіндегі электрондық коммерция нарығының даму тенденциялары//
«Тұран» университетінің хабаршысы» ғылыми журналы. - 2023 ж. - № 1(97) -
б.227-239. DOI 10.46914/1562-2959-2023-1-1-227-239.

13. Azatbek Т., Baitenizov D., Zhanbozova A. Structural analysis of
Kazakhstan's e-commerce market participants in the B2C segment//
Economic Series of the Bulletin of L.N. Gumilyov ENU. - 2024. - Vol.3.
- P.176-192. DOI 10.32523/2789-4320-2024-3-176-192.

14. Yeleussizova B.K., Kazbekova S.Zh., Zeinelgabdin A.B. Cross-border
Retail E-Commerce and De Minimis: The Case of Kazakhstan// World Customs
Journal. - 2024. - Vol.18(1). - P.50-67. DOI
10.55596/001c.117024.

15. Yang Y., Su X., Yao S., Tao C. Mapping big data analytics in the
sharing economy// Frontiers in Environmental Science. - 2022. - Vol.10:
1045943. DOI 10.3389/fenvs.2022.1045943.
\end{refs}

\begin{center}
{\bfseries References}
\end{center}

\begin{refs}
1. Duan C. A State-of-the-Art Review of Sharing Economy Business Models
and a Forecast of Future Research Directions for Sustainable
Development: A Bibliometric Analysis Approach// Sustainability. - 2023.
- Vol.15(5). - P.4568. DOI 10.3390/su15054568.

2. Faraji M., Seifdar M.H., Amiri B. Sharing economy for sustainability:
A review// Journal of Cleaner Production. - 2024. - Vol.434: 140065.
DOI 10.1016/j.jclepro.2023.140065.

3. Tan T.M., et al. Trust in blockchain-enabled exchanges: Future
directions in platform-based markets// Journal of the Academy of
Marketing Science. - 2023. - Vol.51(4). - P.914-939. DOI
10.1007/s11747-022-00889-0.

4. Xu S., Chen T. Research on Optimal Group-Purchase Threshold and
Pricing Strategy of Community Group Purchase//Mathematics. - 2023. -
Vol.11(24):4951. DOI 10.3390/math11244951.

5. Jiang Y., Wu M., Li X. Cooperation mechanism, product quality, and
consumer segment in group buying // Electronic Commerce Research and
Applications. - 2024. - Vol.64: 101357. DOI
10.1016/j.elerap.2024.101357.

6. Kopalle P.K., et al. Dynamic pricing: Definition, implications for
managers, and future research directions // Journal of Retailing. -
2023. - Vol.99(4). - P.580--593. DOI 10.1016/j.jretai.2023.11.003.

7. Rohn D., Bican P., Brem A., Kraus S., Clauß T. Digital platform-based
business models. An exploration of critical success factors // Journal
of Engineering and Technology Management. - 2021. - Vol.60(3). DOI
10.1016/j.jengtecman.2021.101625.

8. Herzallah D., Liébana-Cabanillas F., Muñoz-Leiva F. Et al. Past,
present, and future of social commerce: a bibliometric analysis // Qual
Quant. - 2025. - Vol.59. - P.3079--3111. DOI
10.1007/s11135-025-02103-z.

9. Xu S, Chen T. Research on Optimal Group-Purchase Threshold and
Pricing Strategy of Community Group Purchase //Mathematics.
-2023.-Vol.11(24):4951. DOI 10.3390/math11244951.

10. Wang J., et al. Group buying with consumer disappointment at failed
deals // Transportation Research Part E. - 2025. - Vol.194: 103903. DOI
10.1016/j.tre.2024.103903.

11. Bureau of National Statistics of the Republic of Kazakhstan.
E-commerce in the Republic of Kazakhstan (2024). {[}Jelektronnyj
resurs{]}.2025. -URL:
https://stat.gov.kz/en/industries/economy/local-market/publications/370122/
- Data obrashhenija: 10.02.2026.

12. Zhanbozova A.B., Azatbek T.A., Dadabaeva D.M., Kashkinbaeva K.
Bөlshek sauda segmentіndegі jelektrondyқ kommercija naryғynyң damu
tendencijalary// «Tұran» universitetіnің habarshysy» ғylymi zhurnaly. -
2023 zh. - № 1(97) - b.227-239. DOI 10.46914/1562-2959-2023-1-1-227-239.
{[}in Kazakh{]}

13. Azatbek Т., Baitenizov D., Zhanbozova A. Structural analysis of
Kazakhstan's e-commerce market participants in the B2C segment //
Economic Series of the Bulletin of L.N. Gumilyov ENU. - 2024. - Vol.3.
- P.176-192. DOI 10.32523/2789-4320-2024-3-176-192.

14. Yeleussizova B.K., Kazbekova S.Zh., Zeinelgabdin A.B. Cross-border
Retail E-Commerce and De Minimis: The Case of Kazakhstan // World
Customs Journal. - 2024. - Vol.18(1). - P.50-67. DOI:
10.55596/001c.117024.

15. Yang Y., Su X., Yao S., Tao C. Mapping big data analytics in the
sharing economy // Frontiers in Environmental Science. - 2022. - Vol.
10: 1045943. DOI 10.3389/fenvs.2022.1045943.
\end{refs}

\begin{info}
\hspace{1em}\emph{{\bfseries Сведения об авторах}}

Ускенбаева Р.К. - доктор технических наук, профессор, НАО КазНИТУ им.
К.И. Сатпаева, Алматы, Казахстан, e-mail:
r.k.uskenbayeva@satbayev.university;

Абдул Разак - PhD, НАО КазНИТУ им. К.И. Сатпаева, Алматы, Казахстан,
e-mail:
r.abdul@satbayev.university;

Касымова А.Б. - PhD, НАО КазНИТУ им. К.И. Сатпаева, Алматы, Казахстан,
e-mail:
a.kassymova@satbayev.university;

Кальпеева Ж.Б. - PhD, НАО КазНИТУ им. К.И. Сатпаева, Алматы, Казахстан,
e-mail: zh.kalpeyeva@satbayev.university;

Анартаева А. М. - докторант PhD, Astana IT University, Астана,
Казахстан, e-mail: 255777@astanait.edu.kz.

\hspace{1em}\emph{{\bfseries Information about the authors}}

Uskenbayeva R.K. - Doctor of Technical Sciences, Professor, KazNRTU
named after K.I. Satbayev, Almaty, Kazakhstan, e-mail:
r.k.uskenbayeva@satbayev.university;

Abdul Razaque - PhD, KazNRTU named after K.I. Satbayev, Almaty,
Kazakhstan, e-mail:
r.abdul@satbayev.university;

Kassymova A.B. - PhD, KazNRTU named after K.I. Satbayev, Almaty,
Kazakhstan, e-mail:
a.kassymova@satbayev.university;

Kalpeeva Zh.B. - PhD, KazNRTU named after K.I. Satbayev, Almaty,
Kazakhstan, e-mail: zh.kalpeyeva@satbayev.university;

Anartayeva A.M. - PhD student, Astana IT University, Astana, Kazakhstan,
e-mail: 255777@astanait.edu.kz.
\end{info}
